\lyseisE{Επανάληψη κεφαλαίου Γ1}
\solr{Γ1}{1} α) Ποσοτικό διακριτό. β) Κατηγορικό. γ) Ποσοτικό συνεχές. δ) Ποσοτικό συνεχές. ε) Ποσοτικό διακριτό. στ) Κατηγορικό.
\solr{Γ1}{2} Πληθυσμός: οι $300$ μαθητές του σχολείου. Δείγμα: τα $40$ παιδιά. Μεταβλητή: το είδος μουσικής. Τα δεδομένα είναι κατηγορικά. Το δείγμα είναι αντιπροσωπευτικό, γιατί επιλέχθηκε με κλήρωση από όλες τις τάξεις.
\solr{Γ1}{3} Όχι. Στην ουρά βρίσκονται μόνο παιδιά που αγοράζουν κάτι. Όσα δεν ξοδεύουν τίποτα στο κυλικείο δεν ρωτιούνται, άρα το ποσό που θα βγει θα είναι μεγαλύτερο από το πραγματικό για όλο το σχολείο.
\solr{Γ1}{4} Μήλο $6$ ($30\%$), βερίκοκο $5$ ($25\%$), πορτοκάλι $4$ ($20\%$), φράουλα $3$ ($15\%$), κεράσι $2$ ($10\%$). Σύνολο $20$ και $100\%$. Το ραβδόγραμμα έχει ύψη $6$, $5$, $4$, $3$, $2$ (ελεύθερη κατασκευή).
\solr{Γ1}{5} $130-140$: $2$ ($10\%$). $140-150$: $4$ ($20\%$). $150-160$: $5$ ($25\%$). $160-170$: $4$ ($20\%$). $170-180$: $3$ ($15\%$). $180-190$: $2$ ($10\%$). Σύνολο $20$ και $100\%$. Οι τιμές $140$, $150$, $160$, $170$ μπαίνουν στην επόμενη κλάση.
\solr{Γ1}{6} Όλες οι κλάσεις έχουν πλάτος $10$. Τα κέντρα είναι $135$, $145$, $155$, $165$, $175$, $185$.
\solr{Γ1}{7} $\frac{4}{40}=10\%$. $30\%$ του $40$ είναι $12$. $\frac{14}{40}=35\%$. $20\%$ του $40$ είναι $8$. $\frac{2}{40}=5\%$. Σύνολο: $4+12+14+8+2=40$ και $10+30+35+20+5=100\%$.
\solr{Γ1}{8} Τουλάχιστον $160$ cm πήδηξαν $4+3+2=9$ παιδιά. Λιγότερο από $150$ cm πήδηξαν $2+4=6$ παιδιά, δηλαδή $\frac{6}{20}=30\%$.
\solr{Γ1}{9} Κάθε παιδί αντιστοιχεί σε $\frac{360^\circ}{24}=15^\circ$. Κρυφτό: $8\cdot15^\circ=120^\circ$. Μπάλα: $6\cdot15^\circ=90^\circ$. Σχοινάκι: $90^\circ$. Κυνηγητό: $4\cdot15^\circ=60^\circ$. Άθροισμα $360^\circ$.
\solr{Γ1}{10} Κάθε μαθητής αντιστοιχεί σε $\frac{360^\circ}{72}=5^\circ$. Βόλεϊ: $100:5=20$. Μπάσκετ: $80:5=16$. Τένις: $60:5=12$. Ποδόσφαιρο: $120:5=24$. Έλεγχος: $20+16+12+24=72$.
\solr{Γ1}{11} Σχετικές συχνότητες: $\frac{2}{25}=8\%$, $\frac{8}{25}=32\%$, $\frac{10}{25}=40\%$, $\frac{5}{25}=20\%$. Στο ιστόγραμμα συχνοτήτων τα κολλητά ορθογώνια έχουν βάση $4$ και ύψη $2$, $8$, $10$, $5$. Στο ιστόγραμμα σχετικών συχνοτήτων τα ύψη είναι $8$, $32$, $40$, $20$ (ελεύθερη κατασκευή).
\solr{Γ1}{12} α) $5+12=17$ παιδιά, δηλαδή $\frac{17}{50}=34\%$. β) $10+5=15$ παιδιά. γ) Η κλάση $40-60$ ($18$ παιδιά).
\solr{Γ1}{13} α) Ιστόγραμμα (συνεχή δεδομένα σε κλάσεις). β) Κυκλικό διάγραμμα (ποσοστά ενός συνόλου, λίγες κατηγορίες). γ) Ραβδόγραμμα (κατηγορικά δεδομένα, πολλές κατηγορίες).
\solr{Γ1}{14} Το άθροισμα είναι $20+30+25+15=90\%$, ενώ πρέπει να είναι $100\%$. Αν το λάθος είναι στην τρίτη κατηγορία, αυτή είναι $100-(20+30+15)=35\%$. Οι γωνίες είναι: $72^\circ$, $108^\circ$, $126^\circ$, $54^\circ$ (άθροισμα $360^\circ$).
\solr{Γ1}{15} α) Λάθος. β) Σωστό. γ) Σωστό. δ) Σωστό. ε) Λάθος (είναι κατηγορικό). στ) Σωστό. ζ) Σωστό. η) Λάθος (ανήκει στην επόμενη κλάση).
\solr{Γ1}{16} α) Σωστό. β) Λάθος. γ) Σωστό. δ) Σωστό. ε) Λάθος (ταιριάζει σε λίγες κατηγορίες). στ) Λάθος (έχουν κενά). ζ) Σωστό. η) Λάθος (το άθροισμα είναι $360^\circ$).
